% Generated draft: requires mathematical and source review.
\documentclass[11pt]{article}
\usepackage[margin=25mm]{geometry}
\usepackage{fontspec}
\setmainfont{Latin Modern Roman}
\usepackage{amsmath,amssymb,mathtools,unicode-math}
\setmathfont{Latin Modern Math}
\usepackage{enumitem,needspace,xurl,hyperref}
\setlength{\parindent}{0pt}
\setlength{\parskip}{5pt}
\setlength{\emergencystretch}{3em}
\newcommand{\N}{\mathbb{N}}
\newcommand{\Z}{\mathbb{Z}}
\newcommand{\Q}{\mathbb{Q}}
\newcommand{\R}{\mathbb{R}}
\newcommand{\C}{\mathbb{C}}
\newcommand{\F}{\mathbb{F}}
\newcommand{\A}{\mathbb{A}}
\newcommand{\PP}{\mathbb{P}}
\newcommand{\E}{\mathbb{E}}
\newcommand{\eps}{\varepsilon}
\newcommand{\dd}{\,\mathrm{d}}
\begin{document}
\section*{UnsolvedMath ID: 30001764}
\textbf{Source problem code:} catalogue entry
\subsection{Definitions and mathematical statement}
\textbf{Source-level notation.} The supplied record denotes by \(R_n(u,c_1,c_2)\) the Kleiman--Piene universal residual polynomials. It does not specify their defining construction, coefficient ring, normalization, admissible monomials, or the meanings and domains of \(u,c_1,c_2\). It also does not provide the proposed linear-system algorithm, its equations, its unknown coefficients, or its input data. Accordingly, this notation identifies the objects named by the source but does not constitute a complete definition.

\textbf{Requested determination.} For every integer \(n\geq 20\), determine the coefficients of \(R_n(u,c_1,c_2)\) that remain undetermined by the proposed linear-system algorithm, and explain the loss of full rank reported for that system.

\textbf{Unresolved qualifications.} A fully specified mathematical problem requires the definitions and normalization of \(R_n\), the coefficient domain, and the exact linear systems for the stated range of \(n\). The record does not specify which notion of full rank is intended, the expected rank, or whether the reported deficiency occurs uniformly for every \(n\geq20\). These details cannot be reconstructed from the supplied fragment. The supplied dated assessment describes the coefficient computation as incomplete and states that later Kleiman--Piene work proves a relevant universality theorem under precise hypotheses; those hypotheses and the theorem's exact relationship to this computation are not supplied.
\subsection{Short English statement}
For integers \(n\geq20\), determine the coefficients left unresolved by the cited algorithm for the Kleiman--Piene universal residual polynomials and explain its reported rank deficiency. The polynomial definitions and linear systems are missing, so this is a source-level target requiring further specification. The supplied registry labels the problem partially solved; that status is not independently verified here.
\subsection{Sources}
\begin{itemize}
\item Supplied local UnsolvedMath registry record, permanent ID \(30001764\), title \emph{Undetermined Coefficients in Universal Residual Polynomials}, source problem code \texttt{catalogue entry}, in the Oberwolfach Reports Open Problems collection. The record gives the proposed year as 2011 and identifies no individual proposer. The polynomials are attributed to Kleiman--Piene.
\item \emph{Real Enumerative Questions in Complex and Tropical Geometry} (2011), Oberwolfach Reports, pp.\ 1117--1175, DOI 10.4171/owr/2011/20. Supplied report reference: \url{https://doi.org/10.4171/owr/2011/20}. The record describes Oberwolfach Reports as documenting mathematical workshops and their open-problem discussions; no problem-specific page or passage is supplied.
\item Supplied literature reference: Kleiman--Piene, \emph{node polynomials for curves on surfaces}, \url{https://doi.org/10.3842/SIGMA.2022.059}.
\item Supplied literature reference: \emph{Kleiman--Piene conjecture for plane curves in P3}, \url{https://arxiv.org/abs/1710.02085}.
\item Literature assessment dated \texttt{2026-08-21}: the supplied curation labels the problem partially solved, describes the fragment as an incomplete computation of node-polynomial coefficients, and states that later Kleiman--Piene work proves the relevant universality theorem under precise hypotheses while the displayed computation for \(n\geq20\) remains incompletely determined. This is supplied dated curation, not an independently verified literature assessment; the record requests independent verification of time-sensitive status and cited claims.
\item The registry labels the statement \texttt{corrected\_verified} and reports that it was checked against the official Oberwolfach report, preserving the hypotheses, notation, and open target. This verification claim is supplied provenance and has not been independently checked here.
\end{itemize}
\end{document}
